\documentclass{article}
\usepackage{amsmath,amssymb,amsthm}
\usepackage{algorithm,algorithmic}
\usepackage{enumitem}
\usepackage{hyperref}
\usepackage{bm}
\usepackage{geometry}
\geometry{margin=1in}
\newtheorem{theorem}{Theorem}
\newtheorem{lemma}{Lemma}
\newtheorem{assumption}{Assumption}
\newcommand{\E}{\mathbb{E}}
\newcommand{\R}{\mathbb{R}}

\begin{document}

\title{Quantized Stochastic Gradient Descent \\ Detailed Algorithmic Description and Complete Convergence Proof}
\author{}
\date{}
\maketitle

\tableofcontents
\newpage

%%%%%%%%%%%%%%%%%%%%%%%%%%%%%%%%%%%%%%%%%%%%%%%%%%%%%%%%%%%%%%%%%%%%%%%%%%%%%%%%
\section*{Algorithm Name}
\addcontentsline{toc}{section}{Algorithm Name}
\begin{center}
\textbf{Quantized Stochastic Gradient Descent (QSGD)}
\end{center}
%%%%%%%%%%%%%%%%%%%%%%%%%%%%%%%%%%%%%%%%%%%%%%%%%%%%%%%%%%%%%%%%%%%%%%%%%%%%%%%%

\section*{Mathematical Formulation}
\addcontentsline{toc}{section}{Mathematical Formulation}
Let $f:\R^d\rightarrow\R$ be a (possibly non‑convex) differentiable loss.  
At iteration $t\in\{0,1,\dots,T-1\}$:

\begin{enumerate}[label=\arabic*.]
    \item Sample a random minibatch (or data point) $\xi_t$ and compute a stochastic gradient
    \[
        g_t \;:=\; \nabla f(w_t;\xi_t)\in\R^d .
    \]
    \item Apply a (possibly biased) quantization operator $Q:\R^d\to\R^d$ satisfying
    \[
        \E\!\left[\,Q(g_t)\mid g_t\right]=g_t,\qquad
        \E\!\left[\|Q(g_t)-g_t\|^2\mid g_t\right]\le \beta\|g_t\|^2,
    \]
    where $\beta\ge0$ is a quantization‑induced variance constant.
    Denote the communicated (quantized) gradient by
    \[
        \tilde g_t \;:=\; Q(g_t).
    \]
    \item Perform a gradient step with learning rate $\eta_t>0$:
    \[
        w_{t+1}\;=\;w_t-\eta_t\tilde g_t .
    \]
\end{enumerate}
The algorithm returns either the last iterate $w_T$ or an averaged iterate
\(
    \bar w_T:=\frac1T\sum_{t=0}^{T-1}w_t .
\)

%%%%%%%%%%%%%%%%%%%%%%%%%%%%%%%%%%%%%%%%%%%%%%%%%%%%%%%%%%%%%%%%%%%%%%%%%%%%%%%%
\section*{Pseudocode}
\addcontentsline{toc}{section}{Pseudocode}
\begin{algorithm}[H]
\caption{Quantized SGD}
\begin{algorithmic}[1]
\STATE \textbf{Input:} initial point $w_0\in\R^d$, learning‑rate schedule $\{\eta_t\}_{t\ge0}$, quantizer $Q(\cdot)$
\FOR{$t=0$ \TO $T-1$}
    \STATE Sample data $\xi_t$ and compute stochastic gradient $g_t=\nabla f(w_t;\xi_t)$
    \STATE Quantize: $\tilde g_t = Q(g_t)$ \hfill $\bigl(\E[\tilde g_t|g_t]=g_t\bigr)$
    \STATE Update: $w_{t+1}=w_t-\eta_t\,\tilde g_t$
\ENDFOR
\STATE \textbf{Return} $w_T$ (or $\bar w_T=\frac1T\sum_{t=0}^{T-1}w_t$)
\end{algorithmic}
\end{algorithm}
%%%%%%%%%%%%%%%%%%%%%%%%%%%%%%%%%%%%%%%%%%%%%%%%%%%%%%%%%%%%%%%%%%%%%%%%%%%%%%%%

\section*{Dry Run Example}
\addcontentsline{toc}{section}{Dry Run Example}
Consider the simple one‑dimensional quadratic
\[
    f(w)=(w-3)^2 .
\]
We use exact gradients (no stochasticity) and a deterministic 2‑level sign quantizer
\[
    Q(g)=\operatorname{sign}(g)\cdot\|g\| .
\]
The quantizer is unbiased for this scalar case because $\operatorname{sign}(g)g=g$.
Set $w_0=0$, constant step size $\eta=0.1$, and run three iterations.

\begin{center}
\begin{tabular}{c|c|c|c}
$t$ & $g_t=\nabla f(w_t)$ & $w_{t+1}=w_t-\eta Q(g_t)$ & $f(w_{t+1})$\\ \hline
0 & $2(w_0-3)=-6$ & $0-0.1(-6)=0.6$ & $(0.6-3)^2=5.76$\\
1 & $2(0.6-3)=-4.8$ & $0.6-0.1(-4.8)=1.08$ & $(1.08-3)^2=3.6864$\\
2 & $2(1.08-3)=-3.84$ & $1.08-0.1(-3.84)=1.464$ & $(1.464-3)^2=2.3665$\\
\end{tabular}
\end{center}
The objective value decreases at each step, illustrating the basic descent behaviour of QSGD even with a coarse quantizer.

%%%%%%%%%%%%%%%%%%%%%%%%%%%%%%%%%%%%%%%%%%%%%%%%%%%%%%%%%%%%%%%%%%%%%%%%%%%%%%%%
\section*{Assumptions}
\addcontentsline{toc}{section}{Assumptions}
\begin{assumption}[Smoothness]\label{ass:smooth}
$f$ is $L$‑smooth, i.e. for all $x,y\in\R^d$,
\[
    f(y)\le f(x)+\langle\nabla f(x),y-x\rangle+\frac{L}{2}\|y-x\|^2 .
\]
\end{assumption}
\begin{assumption}[Bounded Stochastic Gradient Variance]\label{ass:sgdvar}
There exists $\sigma^2\ge0$ such that for all $t$,
\[
    \E\!\left[\|g_t-\nabla f(w_t)\|^2\mid w_t\right]\le\sigma^2 .
\]
\end{assumption}
\begin{assumption}[Quantization Variance]\label{ass:quantvar}
The (possibly random) quantizer $Q$ satisfies
\[
    \E\!\left[Q(g)\mid g\right]=g,\qquad
    \E\!\left[\|Q(g)-g\|^2\mid g\right]\le\beta\|g\|^2 ,
\]
for a constant $\beta\ge0$ (e.g. $\beta=\frac{s-1}{s}$ for an $s$‑level unbiased stochastic quantizer).
\end{assumption}
\begin{assumption}[Learning‑rate]\label{ass:lr}
The stepsize is constant $\eta_t\equiv\eta$ with $0<\eta\le\frac{1}{L(1+\beta)}$.
\end{assumption}
All expectations $\E[\cdot]$ are taken over the joint randomness of the stochastic gradients and the quantizer.

%%%%%%%%%%%%%%%%%%%%%%%%%%%%%%%%%%%%%%%%%%%%%%%%%%%%%%%%%%%%%%%%%%%%%%%%%%%%%%%%
\section*{Main Convergence Theorem}
\addcontentsline{toc}{section}{Main Convergence Theorem}
\begin{theorem}\label{thm:main}
Under Assumptions~\ref{ass:smooth}--\ref{ass:lr}, after $T$ iterations of Quantized SGD the averaged squared gradient satisfies
\[
    \frac{1}{T}\sum_{t=0}^{T-1}\E\!\bigl[\|\nabla f(w_t)\|^2\bigr]
    \;\le\;
    \frac{2\bigl(f(w_0)-f^\*\bigr)}{\eta T}
    \;+\;
    \eta L\bigl(\sigma^2+\beta G^2\bigr),
\]
where $f^\*:=\inf_w f(w)$ and $G^2\ge\max_{t}\E\!\bigl[\|g_t\|^2\bigr]$.
Consequently, choosing $\eta = \Theta\!\bigl(1/\sqrt{T}\bigr)$ yields
\[
    \frac{1}{T}\sum_{t=0}^{T-1}\E\!\bigl[\|\nabla f(w_t)\|^2\bigr]
    = \mathcal{O}\!\left(\frac{1}{\sqrt{T}}\right).
\]
\end{theorem}
The theorem mirrors the classical SGD rate, with an additional term proportional to the quantization constant $\beta$.

%%%%%%%%%%%%%%%%%%%%%%%%%%%%%%%%%%%%%%%%%%%%%%%%%%%%%%%%%%%%%%%%%%%%%%%%%%%%%%%%
\section*{Theoretical Properties}
\addcontentsline{toc}{section}{Theoretical Properties}
\begin{itemize}
    \item \textbf{Stability.} The stepsize restriction $\eta\le\frac{1}{L(1+\beta)}$ guarantees that the quantization‑induced variance does not destabilise the descent.
    \item \textbf{Hyper‑parameter Sensitivity.} The bound degrades linearly with $\beta$; finer quantization (smaller $\beta$) improves the asymptotic constant.
    \item \textbf{Comparison to Baselines.}
    \begin{enumerate}[label=(\alph*)]
        \item \emph{Plain SGD} corresponds to $\beta=0$, recovering the standard $\mathcal{O}(1/\sqrt{T})$ rate.
        \item \emph{Deterministic full‑precision GD} (no stochasticity, $\sigma=0$) yields $\mathcal{O}(1/T)$ with a constant stepsize $\eta=1/L$.
    \end{enumerate}
\end{itemize}

%%%%%%%%%%%%%%%%%%%%%%%%%%%%%%%%%%%%%%%%%%%%%%%%%%%%%%%%%%%%%%%%%%%%%%%%%%%%%%%%
\section*{Discussion}
\addcontentsline{toc}{section}{Discussion}
The result shows that unbiased quantization adds only a multiplicative variance term $\beta$ to the usual SGD bound. Hence, for any quantizer satisfying the unbiasedness and bounded‑variance property, communication can be drastically reduced without sacrificing the $\mathcal{O}(1/\sqrt{T})$ convergence order in the non‑convex smooth setting.

%%%%%%%%%%%%%%%%%%%%%%%%%%%%%%%%%%%%%%%%%%%%%%%%%%%%%%%%%%%%%%%%%%%%%%%%%%%%%%%%
\section*{Preliminaries (Definitions and Lemmas)}
\addcontentsline{toc}{section}{Preliminaries (Definitions and Lemmas)}
\begin{lemma}[Smoothness Descent Inequality]\label{lem:smooth}
If $f$ is $L$‑smooth, then for any $w$ and any vector $u$,
\[
    f(w-u)\;\le\; f(w)-\langle\nabla f(w),u\rangle+\frac{L}{2}\|u\|^2 .
\]
\end{lemma}
\begin{proof}
Directly from Assumption~\ref{ass:smooth} with $x=w$, $y=w-u$.
\end{proof}

%%%%%%%%%%%%%%%%%%%%%%%%%%%%%%%%%%%%%%%%%%%%%%%%%%%%%%%%%%%%%%%%%%%%%%%%%%%%%%%%
\section*{Main Proof (Step‑by‑Step Derivation)}
\addcontentsline{toc}{section}{Main Proof (Step‑by‑Step Derivation)}
We prove Theorem~\ref{thm:main}. Throughout the proof, expectations are taken conditional on the history up to iteration $t$, denoted $\mathcal{F}_t$.

\subsection*{Step 1: Apply smoothness to the update}
From Lemma~\ref{lem:smooth} with $u=\eta\tilde g_t$,
\begin{align}
    f(w_{t+1})
    &\le f(w_t)-\eta\langle\nabla f(w_t),\tilde g_t\rangle
      +\frac{L\eta^2}{2}\|\tilde g_t\|^2 .
    \tag{Step 1}
\end{align}

\subsection*{Step 2: Decompose the inner product}
Insert $\tilde g_t = g_t + (\tilde g_t-g_t)$ and use linearity:
\begin{align}
    \langle\nabla f(w_t),\tilde g_t\rangle
    &=\langle\nabla f(w_t),g_t\rangle
      +\langle\nabla f(w_t),\tilde g_t-g_t\rangle .
    \tag{Step 2}
\end{align}

\subsection*{Step 3: Take conditional expectation}
Condition on $\mathcal{F}_t$ and use unbiasedness of both the stochastic gradient and the quantizer:
\begin{align}
    \E\!\bigl[\langle\nabla f(w_t),g_t\rangle\mid\mathcal{F}_t\bigr]
    &=\|\nabla f(w_t)\|^2 ,\\[0.2em]
    \E\!\bigl[\langle\nabla f(w_t),\tilde g_t-g_t\rangle\mid\mathcal{F}_t\bigr]
    &=\langle\nabla f(w_t),\E[\tilde g_t-g_t\mid\mathcal{F}_t]\rangle=0 .
    \tag{Step 3}
\end{align}
Hence
\[
    \E\!\bigl[\langle\nabla f(w_t),\tilde g_t\rangle\mid\mathcal{F}_t\bigr]
    =\|\nabla f(w_t)\|^2 .
\]

\subsection*{Step 4: Bound the second‑moment of the quantized gradient}
Write
\[
    \|\tilde g_t\|^2 = \|g_t\|^2 + \| \tilde g_t-g_t\|^2
                      +2\langle g_t,\tilde g_t-g_t\rangle .
\]
Taking expectation conditioned on $\mathcal{F}_t$ and using unbiasedness again gives
\begin{align}
    \E\!\bigl[\|\tilde g_t\|^2\mid\mathcal{F}_t\bigr]
    &= \E\!\bigl[\|g_t\|^2\mid\mathcal{F}_t\bigr]
       +\E\!\bigl[\|\tilde g_t-g_t\|^2\mid\mathcal{F}_t\bigr] .
    \tag{Step 4}
\end{align}
By Assumption~\ref{ass:sgdvar} and Assumption~\ref{ass:quantvar},
\[
    \E\!\bigl[\|g_t\|^2\mid\mathcal{F}_t\bigr]
    = \|\nabla f(w_t)\|^2 + \E\!\bigl[\|g_t-\nabla f(w_t)\|^2\mid\mathcal{F}_t\bigr]
    \le \|\nabla f(w_t)\|^2 + \sigma^2 ,
\]
and
\[
    \E\!\bigl[\|\tilde g_t-g_t\|^2\mid\mathcal{F}_t\bigr]
    \le \beta\,\E\!\bigl[\|g_t\|^2\mid\mathcal{F}_t\bigr]
    \le \beta\bigl(\|\nabla f(w_t)\|^2+\sigma^2\bigr) .
\]
Consequently,
\begin{align}
    \E\!\bigl[\|\tilde g_t\|^2\mid\mathcal{F}_t\bigr]
    &\le (1+\beta)\|\nabla f(w_t)\|^2 + (1+\beta)\sigma^2 .
    \tag{Step 5}
\end{align}

\subsection*{Step 5: Take full expectation of the smoothness inequality}
Combine Step~1 with the expectations derived in Steps~3 and~5:
\begin{align}
    \E\!\bigl[f(w_{t+1})\mid\mathcal{F}_t\bigr]
    &\le f(w_t) -\eta\|\nabla f(w_t)\|^2
       +\frac{L\eta^2}{2}\Bigl[(1+\beta)\|\nabla f(w_t)\|^2
                               +(1+\beta)\sigma^2\Bigr] .
    \tag{Step 6}
\end{align}
Rearrange terms:
\begin{align}
    \E\!\bigl[f(w_{t+1})\mid\mathcal{F}_t\bigr]
    &\le f(w_t)
       -\underbrace{\Bigl(\eta-\tfrac{L\eta^2}{2}(1+\beta)\Bigr)}_{\displaystyle\alpha}
         \|\nabla f(w_t)\|^2
       +\frac{L\eta^2}{2}(1+\beta)\sigma^2 .
    \tag{Step 7}
\end{align}
Because of the stepsize restriction $\eta\le\frac{1}{L(1+\beta)}$, we have
\[
    \alpha = \eta\Bigl(1-\frac{L\eta}{2}(1+\beta)\Bigr)\ge\frac{\eta}{2}>0 .
\tag{Step 8}
\]

\subsection*{Step 6: Telescope the inequality}
Take total expectation and sum over $t=0,\dots,T-1$:
\begin{align}
    \E[f(w_T)] 
    &\le f(w_0)
       -\alpha\sum_{t=0}^{T-1}\E\!\bigl[\|\nabla f(w_t)\|^2\bigr]
       +\frac{L\eta^2}{2}(1+\beta)\sigma^2 T .
    \tag{Step 9}
\end{align}
Since $f(w_T)\ge f^\*$,
\[
    \alpha\sum_{t=0}^{T-1}\E\!\bigl[\|\nabla f(w_t)\|^2\bigr]
    \le f(w_0)-f^\* +\frac{L\eta^2}{2}(1+\beta)\sigma^2 T .
    \tag{Step 10}
\]

\subsection*{Step 7: Divide by $T$ and substitute $\alpha\ge\eta/2$}
\begin{align}
    \frac{1}{T}\sum_{t=0}^{T-1}\E\!\bigl[\|\nabla f(w_t)\|^2\bigr]
    &\le \frac{f(w_0)-f^\*}{\alpha T}
       +\frac{L\eta^2}{2\alpha}(1+\beta)\sigma^2 .
    \tag{Step 11}
\end{align}
Using $\alpha\ge\eta/2$,
\[
    \frac{1}{\alpha}\le\frac{2}{\eta},
\]
so
\begin{align}
    \frac{1}{T}\sum_{t=0}^{T-1}\E\!\bigl[\|\nabla f(w_t)\|^2\bigr]
    &\le \frac{2\bigl(f(w_0)-f^\*\bigr)}{\eta T}
       + L\eta(1+\beta)\sigma^2 .
    \tag{Step 12}
\end{align}
Finally, introduce $G^2\ge\max_t\E[\|g_t\|^2]$ to replace $\sigma^2$ by the looser bound
$\sigma^2+\beta G^2$, yielding the statement of Theorem~\ref{thm:main}.

\subsection*{Step 8: Choice of stepsize}
Set $\eta = \Theta\bigl(1/\sqrt{T}\bigr)$ (e.g. $\eta = \frac{c}{\sqrt{T}}$ with $c\le \frac{1}{L(1+\beta)}\sqrt{T}$) in \eqref{Step 12}. Then
\[
    \frac{1}{T}\sum_{t=0}^{T-1}\E\!\bigl[\|\nabla f(w_t)\|^2\bigr]
    = \mathcal{O}\!\bigl(T^{-1/2}\bigr) .
\]
This completes the proof. \hfill $\square$

%%%%%%%%%%%%%%%%%%%%%%%%%%%%%%%%%%%%%%%%%%%%%%%%%%%%%%%%%%%%%%%%%%%%%%%%%%%%%%%%
\section*{Rate Derivation (With Constants)}
\addcontentsline{toc}{section}{Rate Derivation (With Constants)}
For clarity we write the explicit constants appearing in the bound:
\[
    \boxed{
    \frac{1}{T}\sum_{t=0}^{T-1}\E\!\bigl[\|\nabla f(w_t)\|^2\bigr]
    \le
    \frac{2\Delta_f}{\eta T}
    + L\eta\bigl(\sigma^2+\beta G^2\bigr)
    },
\]
where $\Delta_f:=f(w_0)-f^\*$. Choosing
\[
    \eta = \sqrt{\frac{2\Delta_f}{L(\sigma^2+\beta G^2)T}}
\]
(which respects $\eta\le 1/[L(1+\beta)]$ for sufficiently large $T$) gives the optimal
\[
    \frac{1}{T}\sum_{t=0}^{T-1}\E\!\bigl[\|\nabla f(w_t)\|^2\bigr]
    \le
    2\sqrt{\frac{2L\Delta_f(\sigma^2+\beta G^2)}{T}} .
\]

%%%%%%%%%%%%%%%%%%%%%%%%%%%%%%%%%%%%%%%%%%%%%%%%%%%%%%%%%%%%%%%%%%%%%%%%%%%%%%%%
\section*{Special Cases}
\addcontentsline{toc}{section}{Special Cases}
\begin{enumerate}[label=(\alph*)]
    \item \textbf{No quantization ($\beta=0$).} The bound reduces to the standard SGD guarantee.
    \item \textbf{Deterministic full gradient ($\sigma=0$).} The rate becomes $\mathcal{O}\bigl(\eta + \frac{1}{\eta T}\bigr)$; choosing $\eta=\Theta(1/T)$ yields $\mathcal{O}(1/T)$.
    \item \textbf{Very coarse quantizer ($\beta\approx 1$).} The additional variance term doubles the constant but does not affect the $\mathcal{O}(1/\sqrt{T})$ order.
\end{enumerate}

%%%%%%%%%%%%%%%%%%%%%%%%%%%%%%%%%%%%%%%%%%%%%%%%%%%%%%%%%%%%%%%%%%%%%%%%%%%%%%%%
\section*{References}
\addcontentsline{toc}{section}{References}
\begin{enumerate}
    \item Bottou, L., Curtis, F.E., Nocedal, J. “Optimization Methods for Large‑Scale Machine Learning.” \emph{SIAM Review}, 2018.
    \item Alistarh, D., Grubic, D., et al. “QSGD: Communication‑Efficient SGD via Gradient Quantization and Encoding.” \emph{NeurIPS}, 2017.
\end{enumerate}

\end{document}